%% Six complete proofs for the declared native common Q-Res presentation.
\begin{annexitem}{annex:p12-qres-game-strategies}{thm:p12-qres-game-strategies}{Finite game semantics and strategy construction}
\begin{proof}
Expand blocks into a fixed order of their variables.  Equal-type quantifiers
within one block commute: existential quantification chooses one Boolean
vector and universal quantification considers every such vector.  The scalar
game tree therefore has the source block game's meaning.  Label each leaf
by the Boolean value of its matrix.  Working upward, an existential node is
true if one child is true; a universal node is true if both children are true.
These are exactly the recursive quantifier clauses, including a zero-depth
tree whose value is that of the quantifier-free matrix.

Construct a winner's strategy simultaneously.  At a true existential node
choose a true child, and at a true universal node retain the existential
strategies of both children, selected by the universal choice.  At a false
universal node choose a false child, and at a false existential node retain
the universal strategies of both children, selected by the existential
choice.  Induction on tree height proves the resulting history strategy
wins against every opponent continuation.  It supplies a winner at every
node, not only along one successful branch.

Its permitted dependencies can be written as the stated functions.  Given
an earlier universal-value vector for an existential variable, reconstruct
the existential player's previous choices using its already defined
functions in prefix order.  This determines the entire earlier history
compatible with that opponent vector, and the chosen history-strategy value
defines \(f_e\).  It is a function of only those earlier universal values.
The same induction reconstructs the universal player's own earlier choices
from earlier existential values and defines \(g_u\).  Any unused history
entry can be filled arbitrarily to make each finite table total; no value
of a later opponent variable is read.  Empty argument lists produce constant
tables, and unused prefix variables cause no exception.

Finally both players cannot have winning legal strategies.  Playing them
against one another recursively gives one well-defined complete play,
since every function reads only earlier values.  Its matrix would have to
be both true and false.  Thus exactly one player wins, and the upward
construction proves truth/falsity has precisely the claimed strategy meaning.
\end{proof}
\end{annexitem}

\begin{annexitem}{annex:p12-qres-rule-soundness}{thm:p12-qres-rule-soundness}{Rule and refutation soundness with the full prefix}
\begin{proof}
Fix any winning existential strategy \(S\) for the original input.  Its
completion of each total universal assignment satisfies every input clause,
so cl-init preserves the invariant.  Tautological input clauses, if any,
are true in every such completion and may be ignored.

For res, fix a total universal assignment and its completed valuation under
\(S\).  If the pivot literal \(p\) is false, the true parent
\(C_1\cup\{p\}\) supplies a true literal of \(C_1\).  If \(p\)
is true, the true parent \(C_2\cup\{\bar p\}\) supplies a true literal
of \(C_2\).  In either case \(C_1\cup C_2\) is true.  This works
also for any pivot-retaining set decomposition permitted by the printed
guards; the completeness construction uses the canonical removed-pivot
decomposition.  The existential-pivot and non-tautology guards still define
which steps this particular calculus accepts.

For red, let the parent be \(C\cup\{l\}\) with universal variable \(u\).
If \(l\) is already in \(C\), the conclusion equals the parent and is
an identity.  Otherwise non-tautology also excludes \(\bar l\) from
\(C\).  Suppose some universal assignment completed by \(S\) falsifies
\(C\).  Change only \(u\)'s value to falsify \(l\).  Every existential
literal in \(C\) precedes \(u\), so its function reads no \(u\).
Its value is unchanged.  Every other universal literal in \(C\) is also
unchanged, and the parent has no opposite-sign occurrence of \(u\) that
could become true.  Thus the changed completion falsifies the parent,
contradicting the invariant.  Consequently every completion satisfies
\(C\).  With no existential literal the dependency argument is vacuous
but the other-literal and non-tautology argument is still needed.  The proof
does not replace the existential strategy or permit it to read a future
universal value.

Induction on the finite derivation establishes the invariant for every
record.  No strategy satisfies an empty clause; hence an empty-clause
derivation excludes every winning existential strategy.  The finite-game
theorem then says the original QBF is false.

The specified checker makes only finite prefix/type, set, membership and
earlier-index tests.  At each record these tests establish exactly an
instance of one of the displayed rules, so the same induction applies to
acceptance.  On an explicit finite encoding, naive token/set comparison and
scanning the prefix and previous records use a polynomial number of
operations on strings of at most the input-plus-certificate length.  The
procedure therefore terminates in polynomial time in that explicit length.
This is a mathematical specification and bound, not a claim that an
unexamined parser or checker binary implements it.
\end{proof}
\end{annexitem}

\begin{annexitem}{annex:p12-qres-elimination-step}{thm:p12-qres-elimination-step}{Exact innermost existential and universal elimination}
\begin{proof}
All current clauses are non-tautological.  Fix an arbitrary complete valuation
\(v\) of the earlier variables.  For innermost existential \(x\), write
the matrix as the conjunction of the no-\(x\) clauses and the two families
\(A\cup\{x\}\), \(B\cup\{\bar x\}\), with the designated pivot
removed from \(A,B\).  A satisfying extension to \(x\) satisfies every
no-\(x\) clause and every pair resolvent \(A\cup B\), by ordinary
resolution.  It thus satisfies the replacement.

Conversely suppose \(v\) satisfies the replacement.  A positive-pivot
clause with all of \(A\) false requires \(x=1\); a negative-pivot clause
with all of \(B\) false requires \(x=0\).  Both demands cannot occur.
If they did, the corresponding \(A\cup B\) would be false under \(v\).
It is then necessarily non-tautological, since no consistent valuation
falsifies a complementary literal pair.  That resolvent was included in
the replacement, a contradiction.  Choose the demanded value if one family
demands it, and either value if neither does.  Every remaining positive
or negative parent already has a true remainder literal.  This extension
satisfies all old clauses.  Empty families, empty remainders, and an already
empty no-pivot clause are covered by the same argument.  Tautological
resolvents were omitted because they are true under every valuation, not
because they could be used as reduction parents.

For innermost universal \(u\), a clause not mentioning \(u\) is unchanged.
A clause \(D\cup\{u\}\), where \(D\) contains neither polarity of
\(u\), is true for both \(u=0\) and \(u=1\) precisely when \(D\)
is true.  The identical statement holds for \(D\cup\{\bar u\}\).
Universal quantification distributes over the finite conjunction of these
clauses, so deleting the universal literal from each computes exactly
\(\forall u\) of the old matrix.  A tautological parent containing both
signs would not meet this argument and has already been removed.  A clause
with just that universal literal becomes the empty clause, correctly
showing falsity.  An unused \(u\) simply disappears from the prefix.

Both equivalences hold for every earlier valuation, not merely one chosen
play.  Applying the unchanged earlier quantifiers therefore preserves truth.
The removed variable is innermost among those still occurring in the
working prefix; every existential literal in a universal-reduction parent
is earlier in the original prefix too.  Thus the semantic replacement
records only legal original-prefix rule steps.
\end{proof}
\end{annexitem}

\begin{annexitem}{annex:p12-qres-completeness}{thm:p12-qres-completeness}{Constructive refutational completeness}
\begin{proof}
Start with the non-tautological input clauses, recording their axiom IDs.
The tautological omissions preserve the matrix under every valuation.
Induct over the variables eliminated from the right.  Maintain that the
working matrix mentions only remaining variables, every working clause has
an earlier valid proof ID under the immutable original prefix, and the
remaining-prefix truth equals the original QBF truth.

For a universal stage, each containing working clause has a legal red step
because every still-present existential is earlier.  Record the reduced
clause and carry no-universal clauses with their existing IDs.  For an
existential stage, record all non-tautological positive/negative pair
resolvents with their earlier parent IDs; the pivot is existential in the
original prefix and the canonical decomposition satisfies the two pivot
guards.  Clauses not involving the pivot keep their IDs.  Dropping the
pivot-containing parents affects the working matrix, not their availability
as earlier proof records.  The elimination lemma proves truth preservation
in either stage and that no eliminated variable remains.  Hence all three
invariant components hold at the next stage.

There are finitely many variables and every stage forms finitely many pairs,
so the construction terminates.  At an empty prefix every current clause
has no variables and is empty; a matrix there is false exactly when it
contains an empty clause, and true when it has no clauses.  If the original
input was false, the invariant therefore supplies an empty working clause
with a valid proof ID.  A prefix-free false input already has that axiom.
If an empty clause appears earlier, its proof can already be returned.

This constructs a refutation for every false finite PCNF regardless of
alternation or width.  Rule soundness gives the converse.  Pair generation
may grow exponentially and the proof makes no polynomial size bound or
claim about a solver's heuristic run.  The existential replacement used
ordinary resolution only at its legitimate innermost existential stage;
the universal stages and their prefix guards are essential.  It is not a
SAT refutation of the original propositional matrix.
\end{proof}
\end{annexitem}

\begin{annexitem}{annex:p12-qres-strategy-checking}{thm:p12-qres-strategy-checking}{Exact meaning of finite strategy acceptance}
\begin{proof}
The prefix and all table argument lists are finite.  A list containing only
permitted earlier opponent variables and a complete Boolean table defines
a total legal function; an illegal future argument or missing row does not.
There are finitely many assignments to all opponent variables, and matrix
evaluation on each induced full assignment terminates.  Thus the checker
terminates.

Fix a legal existential table family.  Every complete universal-value vector
defines one legal play with that family: in prefix order, its existential
values are computed from already supplied earlier universal values.  Any
play against any universal history strategy produces some such universal
vector.  Conversely each fixed vector is realizable by a universal player
that chooses those constants at its turns.  The enumerated plays are
therefore exactly all possible outcomes against the family.  Requiring the
matrix true on all of them is precisely winning.  With the roles reversed,
the identical argument for a legal universal family enumerates all possible
existential vectors and requires falsity, again precisely winning.  The
argument covers dependent opponent behavior as well as constant opponents;
one cannot evade a losing play by observing that the opponent may use functions.

The game-construction theorem supplies such a family for the unique winner,
so it supplies a finite certificate for either truth value.  With no
opponent variable the enumeration has one empty row, not zero tests.
With an empty prefix it tests just the matrix constant.  Tables may be
exponentially large or require exponentially many play checks in terms
of variable count, so no succinct-strategy complexity or arbitrary
refutation-to-strategy extraction claim follows.
\end{proof}
\end{annexitem}

\begin{annexitem}{annex:p12-qres-normalized-bridge}{thm:p12-qres-normalized-bridge}{The full-reduction convention at its actual standing}
\begin{proof}
Full reduction removes all universal literals having no later existential
literal in their clause.  Removing one does not change the existential
literal set, so the others remain eligible.  There are finitely many
literals; any such sequence of individual legal red steps gives exactly
\(\operatorname{ur}_{\Pi}(C)\).  Rule soundness shows every winning
strategy of the original matrix satisfies its fully reduced clauses.
Conversely those reduced clauses are literal subsets of their parents,
so satisfying them also satisfies the parents.  The winning existential
strategy sets are equal, establishing the input-normalization obligation.

A normalized 2013 resolution step first uses the same existential pivot
on its fully reduced parents, rejects a tautological union, and then performs
this finite full reduction.  Expanding it into res followed by those red
steps gives a valid derivation on the common sheet.  Thus normalized
refutation is sound; the initial normalization can itself be recorded by
axioms and red steps against the original input.

For completeness of that convention, start the elimination construction
with the fully reduced input and maintain a fully reduced working matrix.
At an innermost universal variable, it is absent from every such clause:
no remaining existential lies after it, so full reduction has already
removed its occurrences.  That stage is vacuous.  At an existential
elimination stage form every non-tautological pair resolvent and fully
reduce it.  These are exactly normalized resolution steps.  The ordinary
existential replacement lemma gives equivalent truth for the earlier
prefix before this full reduction, and the full-reduction argument just
proved preserves it afterward.  No-pivot clauses already fully reduced
remain so.  Repeating the finite construction yields the empty clause
exactly for a false input, with a normalized proof ID.  Prefix-free falsity
is the same empty axiom case.

This proves both conventions refute precisely false closed finite PCNFs.
It does not say every raw intermediate clause is literally equal to the
fully reduced one; a normalized proof can remove universals sooner than
a chosen common-sheet proof.  Nor does existence of this semantic bridge
identify the 2013, 2021 or unheld 1995 source texts or establish any
unclaimed polynomial proof translation or size relation.
\end{proof}
\end{annexitem}
